\chapter{Introdução ao CSS e Primeiros Passos}

\section{O que é CSS?}

A primeira observação importante é que, mesmo que você nunca escreva uma única linha de CSS, os navegadores web ainda assim conseguem renderizar visualmente qualquer documento HTML. Isso é possível porque os navegadores já implementam, internamente, um conjunto de estilos padrão para os elementos definidos na especificação HTML. É por isso que uma lista não ordenada (ul) já aparece automaticamente com marcadores (bullets) na tela, mesmo sem nenhuma folha de estilo associada: esse é o estilo padrão do navegador. Quando escrevemos nossas próprias regras CSS, na prática estamos sobrescrevendo esses estilos padrão recebidos do navegador. O CSS, portanto, é a linguagem que especifica como os elementos definidos estruturalmente e semanticamente pelo HTML devem aparecer para o usuário final. Com CSS, é possível alterar tipo e tamanho de fonte, cores, espaçamentos, disposição em colunas, posicionamento de elementos na tela e, nas versões mais recentes da linguagem, até criar animações — tarefas que antigamente exigiam linguagens de programação completas.

\begin{infobox}{CSS não é exclusivo do HTML}
Embora, na prática, usemos CSS quase sempre para estilizar documentos HTML, essa não é a única linguagem de marcação que o CSS pode estilizar — existem outras, como SVG e XML. Esse detalhe se tornará relevante mais adiante, quando você trabalhar com programação back-end e entender melhor o conceito de user agent, um termo mais amplo do que ”navegador web”, que se refere a qualquer software capaz de interpretar e representar um documento em nome do usuário.
\end{infobox}

\section{A Sintaxe de uma Regra CSS}

O CSS é uma linguagem baseada em regras. Cada regra é composta por um seletor — que indica qual elemento (ou elementos) do documento HTML deve receber o estilo — e um bloco de declarações, delimitado por chaves, contendo pares de propriedade e valor.

\begin{codebox}{CSS}
 h1 {
   color: red;
   font-size: 5rem;
 }
\end{codebox}

Nesse exemplo, h1 é o seletor: ele indica que a regra se aplica a todos os elementos h1 do documento. Dentro do bloco de declarações, temos duas propriedades: color, cujo valor é red (vermelho), e font-size, cujo valor é 5rem (uma unidade relativa de tamanho que estudaremos com mais detalhe adiante). Se essa regra for aplicada a um documento HTML contendo um elemento h1, o título ficará vermelho e com tamanho de fonte maior.

\begin{infobox}{Dica}
Uma folha de estilo é composta de várias regras, escritas uma após a outra. Sempre que tiver dúvidas sobre quais valores são aceitos por uma determinada propriedade CSS, uma estratégia muito eficaz é pesquisar no Google o nome da propriedade seguido de ”MDN”(por exemplo, ”color MDN”). O primeiro resultado normalmente será a página de referência da Mozilla Developer Network para aquela propriedade, com diversos exemplos práticos de uso.
\end{infobox}

\section{Associando CSS a um Documento HTML}

\begin{codebox}{Código}
Existem três formas de aplicar CSS a um documento HTML: folha de estilo externa,
folha de estilo interna e estilo em linha (inline). Este capítulo apresenta a forma
recomendada — a folha de estilo externa — que será aprofundada, junto das outras
duas, no próximo capítulo.
Para associar uma folha de estilo externa, utiliza-se o elemento link dentro do
head do documento HTML, com dois atributos: rel (indicando o relationship, ou
seja, que o documento associado é uma folha de estilo) e href (indicando o caminho,
relativo ou absoluto, até o arquivo .css).
HTML — index.html
<!DOCTYPE html>
<html lang="pt-BR">
<head>
<meta charset="UTF-8">
<title>Meu primeiro site com CSS</title>
<link rel="stylesheet" href="estilos.css">
\end{codebox}

\begin{codebox}{Código}
</head>
<body>
<h1>Bem-vindo ao meu site</h1>
<p class="destaque">Este parágrafo tem uma classe especial.</p>
<p id="rodape-info">Informações de rodapé.</p>
</body>
</html>
\end{codebox}

CSS — estilos.css h1 \{ color: navy; \}

p \{ color: black; font-size: 1.1rem; \}

Repare que o arquivo estilos.css está localizado no mesmo diretório do arquivo index.html — por isso, o caminho informado em href é apenas o nome do arquivo. Se a folha de estilo estivesse dentro de uma subpasta chamada styles, o caminho seria styles/estilos.css, seguindo exatamente a mesma lógica de caminhos relativos já estudada no capítulo sobre hiperlinks.

\section{Seletores Básicos: Tag, Classe e ID}

Um dos aspectos mais poderosos do CSS é a capacidade de selecionar exatamente quais elementos do documento devem receber um determinado estilo. Os três seletores mais fundamentais são o seletor de tag (elemento), o seletor de classe e o seletor de ID.

\subsection{Seletor de Tag (Elemento)}

O seletor de tag seleciona todos os elementos daquele tipo no documento, simplesmente escrevendo o nome da tag:

\begin{codebox}{CSS}
p {
  color: green;
}
\end{codebox}

Essa regra transforma a cor de todos os parágrafos da página para verde, sem exceção. É possível, ainda, aplicar a mesma regra a múltiplos seletores de uma vez, separando-os por vírgula:

\begin{codebox}{CSS}
p, li {
  color: green;



}
\end{codebox}

Essa regra diz: ”tanto os parágrafos quanto os itens de lista devem ter a fonte na cor verde”.

\subsection{Seletor de Classe}

Quando queremos selecionar apenas um subconjunto de elementos — e não todos os elementos de um determinado tipo — usamos o atributo class no HTML, e o seletor de classe no CSS, precedido de um ponto (.).

\begin{codebox}{HTML}
<ul>
  <li>Item comum</li>
  <li class="especial">Item especial</li>
  <li>Item comum</li>
</ul>
\end{codebox}

\begin{codebox}{CSS}
.especial {
  color: orange;
  font-style: italic;
}
\end{codebox}

Essa regra diz: ”todo elemento que possuir um atributo class com valor especial deve ter a fonte na cor laranja e em itálico”. Note que a mesma classe pode ser reaproveitada em quantos elementos forem necessários, inclusive de tags diferentes — por exemplo, tanto um li quanto um span poderiam compartilhar a classe especial:

\begin{codebox}{CSS}
.especial, span.especial {
  color: orange;
}
\end{codebox}

\subsection{Seletor de ID}

O atributo id identifica um elemento de forma única dentro do documento — ou seja, teoricamente, apenas um único elemento no documento inteiro deve possuir um determinado valor de id. No CSS, o seletor de ID é precedido de uma cerquilha (\#).

\begin{codebox}{HTML}
<p id="rodape-info">Informações de rodapé.</p>
\end{codebox}

\begin{codebox}{CSS}
#rodape-info {
  color: gray;
  font-size: 0.9rem;
}
\end{codebox}

\begin{codebox}{Dica}
Um mesmo elemento pode ter apenas um valor de id, mas pode ter vários val-
ores de class simultaneamente, separados por espaço no HTML (por exemplo,
class="destaque grande"). Além disso, como veremos com mais profundidade
no capítulo sobre cascata e especificidade, um seletor de ID tem precedência sobre
um seletor de classe, que por sua vez tem precedência sobre um seletor de tag —
essa hierarquia de ”peso”entre seletores será fundamental para resolver conflitos
entre regras CSS.
\end{codebox}

\section{Um Exemplo Combinando Tudo}

\begin{codebox}{Código}
Vamos consolidar os conceitos apresentados com um exemplo completo, combinando
um documento HTML com uma folha de estilo externa que utiliza os três tipos de
seletores estudados:
HTML — index.html
<!DOCTYPE html>
<html lang="pt-BR">
<head>
<meta charset="UTF-8">
<title>Exemplo combinado</title>
<link rel="stylesheet" href="estilos.css">
</head>
<body>
<h1>Blog de Programação</h1>
<p>Este é um parágrafo comum.</p>
<p class="destaque">Este parágrafo recebe destaque visual.</p>
<p id="assinatura">Escrito por um futuro desenvolvedor front-end.</p>
</body>
</html>
\end{codebox}

CSS — estilos.css h1 \{ color: navy; \}

p \{ color: \#333333; font-family: Georgia, serif; \}

.destaque \{ background-color: yellow; font-weight: bold; \}

\#assinatura \{

font-style: italic; color: gray; \}

\begin{codebox}{Código}
Nesse exemplo, o título recebe a cor azul-marinho por meio do seletor de tag; todos
os parágrafos recebem uma cor de texto e uma família de fontes específica, também por
seletor de tag; o parágrafo com class="destaque" ganha um fundo amarelo e negrito
adicional, sem afetar os demais parágrafos; e o parágrafo com id="assinatura" recebe
um estilo próprio, exclusivo daquele único elemento.
\end{codebox}

\begin{infobox}{Dica}
Ao testar seus próprios exemplos, você pode trabalhar tanto localmente (com um editor como o Visual Studio Code, por exemplo usando a extensão ”Live Server”para servir a página) quanto em editores online, como o CodePen ou diretamente nos exemplos interativos disponibilizados pela própria documentação da MDN. Ambas as abordagens são válidas para praticar os conceitos apresentados neste capítulo.
\end{infobox}

Síntese do Capítulo

\begin{itemize}
  \item CSS é a linguagem responsável por definir a apresentação visual e o layout de documentos estruturados com HTML, complementando a estrutura e a semântica fornecidas por esta última.
\end{itemize}

\begin{itemize}
  \item Mesmo sem CSS, navegadores já aplicam estilos padrão internos aos elementos HTML; escrever CSS é, na prática, sobrescrever esses estilos.
\end{itemize}

\begin{itemize}
  \item Uma regra CSS é composta por um seletor e um bloco de declarações entre chaves, com pares de propriedade e valor separados por dois-pontos.
\end{itemize}

\begin{itemize}
  \item A forma recomendada de associar CSS a HTML é a folha de estilo externa, vinculada com <link rel="stylesheet"href="...» dentro do head.
\end{itemize}

\begin{itemize}
  \item O seletor de tag afeta todos os elementos daquele tipo; o seletor de classe (precedido de ponto) afeta um subconjunto de elementos marcados com class; o seletor de ID (precedido de cerquilha) afeta um único elemento identificado por id.
\end{itemize}

\begin{itemize}
  \item Seletores podem ser combinados em uma mesma regra separando-os por vírgula, aplicando o mesmo bloco de declarações a múltiplos alvos.
\end{itemize}
